\section{Introduction}\label{sec: introduction}

Let $A^\bullet=\bigoplus_{k=0}^dA^k$ be a finitely generated, commutative graded $\R$-algebra with a degree map $\deg\colon A^d\to \R$. $A^\bullet$ is said to satisfy the \textit{K\"ahler package} with respect to an open convex cone $K\subseteq A^1$ if
\begin{enumerate}
\item (Poincar\'e duality) $\deg$ is an isomorphism, and  for $0\leq k\leq d$, the bilinear pairing\[ A^{k}\times A^{d-k}\to\R\]that sends $(\alpha,\beta)$ to $\deg(\alpha\beta)$ is non-degenerate. In particular $A^0\cong \R$.
\item (mixed Hard Lefschetz) For $0\leq k\leq \lfloor\frac{d}{2}\rfloor$ and any $\ell_1,\ldots,\ell_{d-2k}\in K$, the linear map \[A^k\to A^{d-k},\alpha\mapsto \left(\prod_{i=1}^{d-2k}\ell_i\right)\alpha\]is an isomorphism.
\item (mixed Hodge-Riemann relations) For $0\leq k\leq \lfloor\frac{d}{2}\rfloor$ and any $\ell_0,\ell_1,\ldots,\ell_{d-2k}\in K$, the symmetric bilinear form \[A^k\times A^k\to \R,\quad (\alpha,\beta)\mapsto(-1)^k \deg\left(\left(\prod_{i=1}^{d-2k}\ell_i\right)\alpha\beta\right)\]is positive-definite on the kernel of the map\[A^k\to A^{d-k+1},\alpha\mapsto \left(\prod_{i=0}^{d-2k}\ell_i\right)\alpha.\]
\end{enumerate}
Classes inside $K$ are often called \textit{ample classes}, whereas classes inside $\overline K$ are called \textit{nef classes}. Prototypical examples of such an algebra would be the rings of real $(k,k)$-cohomology classes on a compact K\"ahler manifold.

The K\"ahler package imposes strong constraints on the intersection numbers of nef classes. The \textit{Alexandrov-Fenchel type} inequalities assert that,
for any classes $\alpha\in A^1$ and $\beta,\ell_3,\ldots,\ell_d\in \overline K$, we have \[
    \deg(\alpha\beta\ell_3\ldots\ell_d)^2\geq\deg(\alpha^2\ell_3\ldots\ell_d)\deg(\beta^2\ell_3\ldots\ell_d).
\]
These inequalities follow from Hodge-Riemann relations up to degree one.

By investigating Hodge-Riemann relations up to degree two, we discover a new family of inequalities. Let $R=\R[X_1,\ldots,X_n]$ equipped with the standard grading, and $K_R=\{\sum_{i=1}^n \lambda_iX_i\mid \lambda_i>0,\forall i\}$, we define partial orders on $R^1$ by saying $A\preceq B$ when $B-A\in \overline K_R$, and $A\prec B$ if $B-A\in K_R$. 
\begin{theorem}\label{thm: main}
Suppose that
\[
    HP=\sum_{i=1}^rA_iB_iV_i\in R\tag{1}\label{eq: referencing}
\]
where $P\in R^2$, $H, A_i,B_i,V_i\in R^1$ and $H,A_i,B_i\in \overline{K_R}$ and $0\preceq A_i,B_i\preceq H$. Let $A^\bullet$ be a graded algebra of degree $d\geq 4$ satisfying the K\"ahler package with respect to $K$. We consider a graded ring homomorphism $f\colon R\to A$ that maps $K_R$ to $\overline K$, and we denote the images of $P,H,A_i,B_i,V_i$ by $p,h,\alpha_i,\beta_i, v_i$ respectively.

For any class $\ell\in A^{d-4}$ given by the product of $(d-4)$ classes in $\overline{K}$, we have the following inequality
\[
\deg (p^2\ell)+25\left(\sum_{i=1}^r \sqrt{\frac{2\deg(\alpha_i\beta_i h v_i\ell)^2}{\deg(\alpha_i\beta_ih^2\ell)}-\deg(\alpha_i\beta_i v_i^2\ell)}\right)^2\geq \frac{2\deg(ph^2\ell)^2}{\deg(h^4\ell)}
\]
whenever $\deg(\alpha_i\beta_ih^2\ell)\neq 0$ for all $i$.
\end{theorem}

To show that these inequalities are not vacuous, we give an example showing that they do not follow from the AF-type inequalities.

\begin{example}\label{ex: ternary}
    Put
\[
        a=\frac{2023}{1500},\qquad
        b=\frac{119}{100},\qquad
        r=\frac{11}{10},\qquad
        s=\frac12,
\]
and define
\[
\begin{aligned}
F(x,y,z)={}&
 a\frac{x^4}{24}
+b\frac{x^3y+x^3z}{6}
+\frac{x^2y^2+x^2z^2+y^2z^2}{4}
+r\frac{x^2yz}{2}                                      \\
&+s\frac{xy^3+xz^3}{6}
+\frac{xy^2z+xyz^2}{2}
+s^2\frac{y^4+z^4}{24}
+s\frac{y^3z+yz^3}{6}.
\end{aligned}
\]
Then for any degree $4$ graded $\R$-algebra $A$ and classes $\chi,\gamma,\zeta\in A^1$ such that
\[
    \frac{1}{4!}\deg(x\chi+y\gamma+z\zeta)^4=F_\epsilon(x,y,z)\coloneqq F(x+\epsilon\sigma,y+\epsilon\sigma,z+\epsilon\sigma),\; \sigma=x+y+z,
\]
the inequalities in \Cref{thm: main} fail for the map $R=\R[X,Y,Z]\to A,\ g(X,Y,Z)\mapsto g(\chi,\gamma,\zeta)$ when $0<\epsilon\ll 1$, where the equality \eqref{eq: referencing} is given by
\begin{align*}
    H  =& X+Y+Z,\\
    P  =& 750X^2-750XY-750XZ-28Y^2-28Z^2+735YZ,\\
    HP =&
        \ X^2\cdot 25(30X-17Y-17Z)\\
        &+XY\cdot \frac{85}{2}(10X-2Y-9Z)+XZ\cdot \frac{85}{2}(10X-9Y-2Z)\tag {1'}\label{eq: referencing'}\\
        &-Y^2\cdot (693(X-Z)+14(2Y-Z))-Z^2\cdot(693(X-Y)+14(2Z-Y)).
\end{align*}
On the other hand, $F_\epsilon(x,y,z)$ is a \textit{strictly Lorentzian polynomial}, which means that it has positive coefficients and the Hessian matrices of its second order partial derivatives have the same signature as Lorentzian metrics. For these polynomials, one can canonically construct Poincar\'e duality algebras satisfying AF-type inequalities.
\end{example}

We organize this paper as follows. In section \ref{sec: norms}, we prove auxiliary results on norm estimates associated to symmetric forms with Lorentzian signature; in section \ref{sec: proof}, we will apply these results to an external inequality \eqref{eq: external inequality} to prove the main theorem; then in section \ref{sec: ex}, we will give detailed analysis on \cref{ex: ternary} and compare different sets of polynomials.
